\subsection{Búsqueda Binaria sobre la Respuesta + Two Pointers (k-ésima suma de pares)}

\graybold{Preguntas guía:}

\begin{itemize}
\item ¿La respuesta es un número y existe una condición monótona (si X funciona, todo valor mayor/menor también funciona)?
\item ¿Puedo verificar en \bigO{n} o \bigO{n \log n} si un valor candidato "cumple" la condición pedida?
\item ¿Dentro de esa verificación necesito contar pares/elementos con una condición de suma sobre un arreglo ordenado? (ahí entra Two Pointers)
\end{itemize}

\graybold{Utilidad:} Cuando no es fácil calcular la respuesta directamente pero sí verificar si un candidato es válido, y esa validez es monótona, se "busca" la respuesta con búsqueda binaria sobre el rango de valores posibles. Two Pointers acelera el conteo en arreglos ordenados de \bigO{n^2} a \bigO{n}.

\graybold{Complejidad:} Búsqueda binaria: \bigO{\log(\text{rango})} iteraciones. Cada verificación con Two Pointers: \bigO{n}. Total: \bigO{n \log(\text{rango})}.

\graybold{Fundamento teórico:} La búsqueda binaria sobre la respuesta funciona porque el conjunto de valores "válidos" forma un intervalo continuo (monotonía), reduciendo el espacio de búsqueda a la mitad en cada paso. Two Pointers explota que, en un arreglo ordenado, mover un extremo cambia la condición de suma de forma predecible, evitando recorrer todas las parejas.

\graybold{Ejercicio aplicado}

Dados N precios de categorías de joyas, un cliente que es el K-ésimo en la fila debe comprar dos joyas de categorías distintas sin repetir una combinación ya usada por otro cliente. Hallar el mínimo dinero que garantiza poder comprar: la K-ésima suma más pequeña entre todos los pares de categorías distintas.

\graybold{I/O:} N no está acotado explícitamente pero puede ser grande → lectura total + tokenización (patrón 2) para el arreglo de precios.

\graybold{Código solución}

\begin{lstlisting}[language=Python]
import sys
 
def count_pairs_leq(prices, n, limit):
    # cuenta cuantos pares cumplen arr[i] + arr[j] <= limit
    # usando dos punteros sobre un arreglo YA ORDENADO
    count = 0
    left, right = 0, n - 1
    while left < right:
        if prices[left] + prices[right] <= limit:
            count += right - left       # todos los pares con "left" fijo
            left += 1
        else:
            right -= 1
    return count
 
def solve():
    data = sys.stdin.buffer.read().split()
    n = int(data[0]); k = int(data[1])
    prices = list(map(int, data[2:2 + n]))
    prices.sort()
 
    low = prices[0] + prices[1]         # menor suma posible
    high = prices[n - 2] + prices[n - 1]  # mayor suma posible
 
    while low < high:
        mid = low + (high - low) // 2
        if count_pairs_leq(prices, n, mid) >= k:
            high = mid                    # mid cumple: intenta algo menor
        else:
            low = mid + 1                # mid no cumple: sube el limite
 
    print(low)
 
solve()
\end{lstlisting}
